\begin{thebibliography}{99}
\bibitem{alpo-announcement}
L. Alp\"oge, announcement of a proof of RH and Dirichlet GRH,
20 November 2026, 4 pp.
\url{https://alpo.ge/remark.pdf}.
The source archive consists of the fourteen parts linked there.

\bibitem{dlmf-dirichlet}
NIST Digital Library of Mathematical Functions,
\S25.15, Dirichlet $L$-functions, especially equations
25.15.1--25.15.5 and the discussion of zeros.
\url{https://dlmf.nist.gov/25.15}.

\bibitem{dlmf-zeta}
NIST Digital Library of Mathematical Functions,
\S25.2, Definition and Expansions of the Riemann zeta function,
especially equations 25.2.1 and 25.2.4.
\url{https://dlmf.nist.gov/25.2}.

\bibitem{dlmf-zeros}
NIST Digital Library of Mathematical Functions,
\S25.10, Zeros of the Riemann zeta function.
\url{https://dlmf.nist.gov/25.10}.
% Insert these items inside the host document's bibliography.
\bibitem{qrh-september}
OpenAI, \emph{The Quasi-Riemann Hypothesis: A Zero-Free Half-Plane
$\operatorname{Re}s>7/8$}, 30 September 2026, 199 pp.
Version at repository commit
\path{4eaee87548250b2682affda7d05d947c05554bae}.
Theorem 1.1; Proposition 2.1; Sections 17--20.
\url{https://github.com/openai/math/blob/4eaee87548250b2682affda7d05d947c05554bae/preprints/The-Quasi-Riemann-Hypothesis-September-30-2026/paper.pdf}.

\bibitem{qrh-october}
OpenAI, \emph{The Quasi-Riemann Hypothesis}, 5 October 2026,
48 pp. Version at the same repository commit.
Theorem 1.1, Proposition 3.1, Section 3, Propositions 5.1, 5.2,
5.4, and Appendix A.2.
\url{https://github.com/openai/math/blob/4eaee87548250b2682affda7d05d947c05554bae/preprints/The-Quasi-Riemann-Hypothesis-October-5-2026/paper2.pdf}.

\bibitem{qrh-catalogue}
H. Collet and contributors, \emph{QRH Catalogue}, repository
commit \path{5b7217781269248f39fbcec2432394d3d1610344},
consulted 20 November 2026. In particular the README's evidence
qualifications and the mathematical entries G276 and G280.
\url{https://github.com/HaileyStorm/QRH-Catalogue/tree/5b7217781269248f39fbcec2432394d3d1610344}.
This reference identifies a catalogue of source claims and
conditional reductions; it is not a certificate for every
listed theorem.

\bibitem{alpo-prime-supply}
L. Alp\"oge, submitted source archive, the ordinary progression
and quadratic-field prime conditions in
\path{Formalism/SharpRestriction/InertJacobiLevels/PrimeSupply.lean};
the varying quadratic ray-class conditions in
\path{Formalism/SharpRestriction/GlobalSchurBounds/ShortIntervalPrimes.lean}.
Archive SHA-256:
\path{9631b4f020199e580ba86178cb56c6e9b09f690c2369a80dc4334aa93f49c7fc}.

\bibitem{alpo-final-scalar}
L. Alp\"oge, same submitted archive,
\path{Formalism/SizeWindowBrackets/GrowthBoundGivesRiemannHypothesis.lean},
in particular the height hypotheses, the conditions on the
three auxiliary characters, and the final size inequalities.
\bibitem{sz-openai}
OpenAI, \emph{Uniform exclusion of Landau--Siegel zeros},
October 1, 2026, 9 pp.
\url{https://github.com/openai/math/blob/4eaee87548250b2682affda7d05d947c05554bae/preprints/Uniform-exclusion-of-Landau-Siegel-zeros-October-1-2026/paper.pdf}.

\bibitem{sz-guo}
Victor Zhenyu Guo, \emph{An explicit elementary proof of no Landau--Siegel zeros},
arXiv:2610.30476v1, October 21, 2026, 17 pp.
\url{https://arxiv.org/html/2610.30476v1}.

\bibitem{sz-tao}
Terence Tao, \emph{A sketch of the proof of non-existence of Siegel zeroes},
What's New, October 19, 2026.
\url{https://terrytao.wordpress.com/2026/10/19/a-sketch-of-the-proof-of-non-existence-of-siegel-zeroes/}.

\bibitem{sz-pescador}
GPT-6 Pro and Claude Opus 5.5, posted and maintained by Jacob Pescador,
\emph{Landau--Siegel zeros and a Stepanov-type nonvanishing lemma},
version 3, October 22, 2026.
\url{https://jacobpescador.github.io/ai-math/papers/landau-siegel-stepanov/}.

\bibitem{sz-vandijk}
Roelof van Dijk,
\emph{A Frobenius determinant proof of the non-existence of Landau--Siegel zeros},
hexagon:2610.00409v2, 6 pp.
\url{https://hexagonmath.org/pdf/2610.00409v2}.

\bibitem{sz-kerr}
Bryce Kerr, \emph{A simpler proof that there are no Landau--Siegel zeros},
arXiv:2610.21099v1, October 15, 2026, 6 pp.
\url{https://arxiv.org/abs/2610.21099}.
\bibitem{ct-fw}
B. Ferrero and L. C. Washington,
The Iwasawa invariant $\mu_p$ vanishes for abelian number fields,
\emph{Annals of Mathematics} (2) \textbf{109} (1979), 377--395.
\url{https://doi.org/10.2307/1971116}.
Primary bibliographic record:
\url{https://annals.math.princeton.edu/1979/190-2/p05}.

\bibitem{ct-sharifi}
R. Sharifi, \emph{Iwasawa Theory}, author's lecture notes,
Chapter 3 (cyclotomic extensions and limits of class groups) and
Theorem 6.2.1 (Ferrero--Washington).
\url{https://www.math.ucla.edu/~sharifi/iwasawa.pdf}.
Author's reading edition, checked 21 November 2026:
\url{https://www.math.ucla.edu/~sharifi/notes/iwasawa-ch03.html},
\url{https://www.math.ucla.edu/~sharifi/notes/iwasawa-ch06.html}.

\bibitem{milne-ant}
J. S. Milne, \emph{Algebraic Number Theory}, version 3.08,
19 July 2020, Chapter 5 (the unit theorem).
\url{https://www.jmilne.org/math/CourseNotes/ANT.pdf}.

\bibitem{milne-cft}
J. S. Milne, \emph{Class Field Theory}, version 4.03,
6 August 2020, Chapter V, \S3, Theorems 3.5--3.6 and
Example 3.9 (Hilbert class field reciprocity).
\url{https://www.jmilne.org/math/CourseNotes/CFT.pdf}.

\bibitem{milne-adt}
J. S. Milne, \emph{Arithmetic Duality Theorems}, second edition,
BookSurge, 2006. Chapter I, Theorem 4.10 and Corollary 4.15,
and Chapter II, Proposition 2.9 give strict global duality, finite
coefficient cohomology and the \'etale--Galois comparison used here.
\url{https://www.jmilne.org/math/Books/ADTnot.pdf}.

\bibitem{weibel-kbook}
C. A. Weibel, \emph{The $K$-book: An Introduction to Algebraic $K$-theory},
Graduate Studies in Mathematics \textbf{145}, American Mathematical
Society, 2013. Author's chapters III, IV and VI, especially III,
Example 1.3.6; IV,
Corollary 1.14.1, Theorems 1.17--1.18 and 6.9, and VI,
equation (8.1.1), Theorem 8.2 and Corollary 8.3.
\url{https://sites.math.rutgers.edu/~weibel/Kbook/Kbook.III.pdf};
\url{https://sites.math.rutgers.edu/~weibel/Kbook/Kbook.IV.pdf};
\url{https://sites.math.rutgers.edu/~weibel/Kbook/Kbook.VI.pdf}.

\bibitem{milnor-moore}
J. W. Milnor and J. C. Moore, On the structure of Hopf algebras,
\emph{Annals of Mathematics} (2) \textbf{81} (1965), 211--264,
especially Corollary 4.18, Theorem 5.18 and the appendix,
pp.\ 262--264.
\url{https://doi.org/10.2307/1970615}.
Author's article in the University of Rochester collection:
\url{https://www.sas.rochester.edu/mth/sites/doug-ravenel/otherpapers/Milnor-Moore.pdf}.

\bibitem{mathew-hurewicz}
A. Mathew, Torsion exponents in stable homotopy and the Hurewicz
homomorphism, \emph{Algebraic \& Geometric Topology}
\textbf{16} (2016), 1025--1041, especially Theorem 6.3.
\url{https://doi.org/10.2140/agt.2016.16.1025}.
Author's article:
\url{https://msp.org/agt/2016/16-2/agt-v16-n2-p12-s.pdf}.

\bibitem{alpo-uniform-source}
L. Alp\"oge, submitted source archive identified in the audit,
especially the directories
\path{Development/AutomorphicArithmetic/AutomorphicVarieties/LayeredIsomorphisms},
\path{Development/AutomorphicArithmetic/AutomorphicVarieties/AdditiveValuationNorms},
\path{Development/AutomorphicArithmetic/AutomorphicVarieties/QuadraticExtensions},
\path{Formalism/LocalToGlobal/Divisibility}, and the theta and
selected-character files listed in the accompanying source manifests.
The source traces specify the exact scope read; these locators do not
assert that every file in the directories was reviewed.

\bibitem{kirschmer-herm}
M. Kirschmer, \emph{Definite quadratic and hermitian forms with
small class number}, habilitation thesis, RWTH Aachen,
\S4.2, especially Theorem 4.2.3 and the definition of the local density,
pp.\ 50--51.
\url{https://www.math.uni-bielefeld.de/~mkirschm/papers/herm.pdf}.
The general mass formula there is attributed to
W. T. Gan and J.-K. Yu, \emph{Group schemes and local densities},
Duke Math.\ J.\ \textbf{105} (2000), 497--524, \S10.

\bibitem{terakado-yu-mass}
Y. Terakado and C.-F. Yu, \emph{Mass formulas and the basic locus
of unitary Shimura varieties}, arXiv:2210.04054v1 (2022),
\S2.3, especially equation (2.6).
\url{https://arxiv.org/pdf/2210.04054}.

\bibitem{harish-chandra-ds-ii}
Harish--Chandra, Discrete series for semisimple Lie groups. II:
Explicit determination of the characters,
\emph{Acta Mathematica} \textbf{116} (1966), 1--111,
especially Theorem 16 and \S41, pp.\ 96--100.
\url{https://doi.org/10.1007/BF02392813}.
Primary article in the Yau Mathematical Sciences Center collection:
\url{https://archive.ymsc.tsinghua.edu.cn/pacm_download/117/6013-11511_2006_Article_BF02392813.pdf}.

\bibitem{duflo-heckman-vergne}
M. Duflo, G. Heckman and M. Vergne,
Projection d'orbites, formule de Kirillov et formule de Blattner,
\emph{M\'emoires de la Soci\'et\'e Math\'ematique de France},
new series, no.\ 15 (1984), 65--128, especially III.1,
theorem (2), p.\ 111.
\url{https://www.numdam.org/article/MSMF_1984_2_15__65_0.pdf}.
This article gives another proof of the formula established by
H. Hecht and W. Schmid, A proof of Blattner's conjecture,
\emph{Inventiones Mathematicae} \textbf{31} (1975), 129--154,
\url{https://doi.org/10.1007/BF01404112}.

\bibitem{baldoni-vergne}
V. Baldoni and M. Vergne,
Discrete series representations and $K$ multiplicities for
$U(p,q)$. User's guide, arXiv:1008.5360v1 (2010),
especially \S2.2, pp.\ 12--14.
\url{https://webusers.imj-prg.fr/~michele.vergne/publications/discreteseries.pdf}.

\bibitem{parthasarathy-minimal}
R. Parthasarathy, Classification of discrete series by minimal $K$-type,
\emph{Representation Theory} \textbf{19} (2015), 167--185,
especially Theorem 1.1, pp.\ 167--168.
\url{https://doi.org/10.1090/ert/467}.
Author-posted full text, used to verify the theorem and its hypotheses:
\url{https://www.researchgate.net/publication/282775275_Classification_of_discrete_series_by_minimal_-type}.
\bibitem{bannai-kobayashi-2007}
K. Bannai and S. Kobayashi,
\emph{Algebraic theta functions and $p$-adic interpolation of
Eisenstein--Kronecker numbers},
arXiv:math/0610163v4, 11 December 2007;
published in Duke Mathematical Journal \textbf{153} (2010), 229--295.
The theorem numbers used here refer to the arXiv version:
Propositions 2.14--2.15, Corollary 2.16(i),
the reduced Mumford-translation formula in Lemma 2.7 and
Proposition 2.8, and \S3.1.
\url{https://arxiv.org/abs/math/0610163v4}.
\url{https://doi.org/10.1215/00127094-2010-024}.
\bibitem{stacks-07S7}
The Stacks Project Authors,
\emph{The Stacks Project}, Lemma 66.14.2, Tag 07S7,
free finite locally free quotients of projective schemes.
\url{https://stacks.math.columbia.edu/tag/07S7}.
\bibitem{sun-eisenstein}
Binyong Sun, \emph{Eisenstein integrals for theta stable parabolic
subalgebras}, arXiv:math/0403440v2, 28 November 2004.
Theorem 6.4 and \S6.4 give the compact-type integral;
\S\S1.1--1.4 fix the group, measure and conjugation conventions.
\url{https://arxiv.org/abs/math/0403440v2}.

\bibitem{vogan-zuckerman}
David A. Vogan, Jr. and Gregg J. Zuckerman,
\emph{Unitary representations with non-zero cohomology},
Compositio Mathematica \textbf{53} (1984), 51--90.
See (5.1), Theorem 5.3 and Proposition 6.1.
\url{https://www.numdam.org/item/CM_1984__53_1_51_0/}.
\bibitem{greither-popescu}
Cornelius Greither and Cristian D. Popescu,
\emph{An equivariant main conjecture in Iwasawa theory and applications},
arXiv:1103.3069v1, 16 March 2011.
Theorem 6.11 and Lemma 6.9, with the conventions in \S6.2.
\url{https://arxiv.org/abs/1103.3069v1}.

\bibitem{getz-hahn-2022}
J. R. Getz and H. Hahn,
\emph{An Introduction to Automorphic Representations},
author manuscript dated 22 April 2022, Theorems 6.2.3
and 6.3.2, with their fixed-level and compact-type conventions.
The theorem numbers here refer to this manuscript.
\url{https://sites.duke.edu/jgetz/files/2022/04/Graduate_Text.pdf}.
\bibitem{weil-siegel}
Andr\'e Weil,
\emph{La formule de Siegel et les groupes classiques},
Acta Mathematica \textbf{113} (1965), 1--87.
The unitary case is \S27, type $(I_2)$; see \S54 for
independence of the normalized quotient volume and
\S57 for the special linear Tamagawa number.
\url{https://archive.ymsc.tsinghua.edu.cn/pacm_download/117/5988-11511_2006_Article_BF02391774.pdf}.

\bibitem{louboutin-residue}
St\'ephane Louboutin,
\emph{Explicit upper bounds for residues of Dedekind zeta functions
and values of $L$-functions at $s=1$, and explicit lower bounds
for relative class numbers of CM-fields},
Canadian Journal of Mathematics \textbf{53} (2001), 1194--1222.
Theorem~1(1).
\url{https://doi.org/10.4153/CJM-2001-045-5}.
\bibitem{banaszczyk-transference}
W. Banaszczyk, New bounds in some transference theorems
in the geometry of numbers,
\emph{Mathematische Annalen} \textbf{296} (1993), 625--635,
especially Theorem~2.2.
\url{https://doi.org/10.1007/BF01445125}.

\bibitem{bost-kunnemann-extensions}
J.-B. Bost and K. K\"unnemann,
Hermitian vector bundles and extension groups on arithmetic
schemes. I. Geometry of numbers,
\emph{Advances in Mathematics} \textbf{223} (2010), 987--1106.
Author preprint, 11 January 2007, especially
Theorem~3.7.2, (3.54), printed p.~46.
\url{https://epub.uni-regensburg.de/13828/1/MP77.pdf}.

\bibitem{khare-larsen-savin}
C. Khare, M. Larsen and G. Savin,
Functoriality and the inverse Galois problem,
\emph{Compositio Mathematica} \textbf{144} (2008), 541--564.
Author version arXiv:math/0610860v3 (29 May 2015),
especially \S4.1, pp.~14--15.
\url{https://arxiv.org/abs/math/0610860}.

\bibitem{zaman-hecke-estimates}
A. Zaman, Explicit estimates for the zeros of Hecke
\(L\)-functions,
\emph{Journal of Number Theory} \textbf{162} (2016), 312--375.
Author version arXiv:1502.05679v4 (29 March 2016),
especially \S2.1 and Lemma~2.1.
\url{https://arxiv.org/abs/1502.05679}.
\bibitem{thorne-adjoint-2022}
J. A. Thorne, \emph{On the vanishing of adjoint Bloch--Kato
Selmer groups of irreducible automorphic Galois representations},
arXiv:2207.04925v1, 11 July 2022, Theorem~6.1.
\url{https://arxiv.org/pdf/2207.04925}.

\bibitem{newton-thorne-adjoint-2023}
J. Newton and J. A. Thorne, \emph{Adjoint Selmer groups of
automorphic Galois representations of unitary type},
J. Eur. Math. Soc. \textbf{25} (2023), 1919--1967.
doi:10.4171/JEMS/1228. Conventions on pp.~1920 and~1923.
\url{https://ems.press/content/serial-article-files/32858}.

\bibitem{alpo-central-comparison}
L. Alp\"oge, submitted archive, the integral family slope comparison and
the divisor inequality, especially
\path{Formalism/RankGapBounds/MeanBoundSecondForm.lean},
\path{Formalism/SizeWindowBrackets/TraceSquareAndDiagBounds.lean},
\path{Formalism/SizeWindowBrackets/ChainSlopeWindows/LevelOfCongruence.lean},
and
\path{Formalism/LocalToGlobal/Divisibility/ConstantGeTwentyEight.lean}.
The independently rewritten comparison in Section~\ref{cs-section}
does not assume the existence of the required arithmetic family.

\bibitem{blggt-potential-automorphy}
T. Barnet-Lamb, T. Gee, D. Geraghty and R. Taylor,
\emph{Potential automorphy and change of weight},
Annals of Mathematics \textbf{179} (2014), 501--609.
Theorem~2.1.1, printed p.~538; conventions on pp.~509--510
and 536--537.
\url{https://doi.org/10.4007/annals.2014.179.2.3}.
\url{https://annals.math.princeton.edu/wp-content/uploads/annals-v179-n2-p03-p.pdf}.

\bibitem{blggt-local-global-p}
T. Barnet-Lamb, T. Gee, D. Geraghty and R. Taylor,
\emph{Local-global compatibility for $l=p$, II},
Annales scientifiques de l'\'Ecole normale sup\'erieure
\textbf{47} (2014), 165--179.
Theorem~1.1, author PDF p.~6, and the monodromy-dominance
definition on pp.~5--6.
\url{https://doi.org/10.24033/asens.2212}.
\url{https://virtualmath1.stanford.edu/~rltaylor/lg2.pdf}.

\bibitem{berger-monodromy}
L. Berger, \emph{Repr\'esentations $p$-adiques et
\'equations diff\'erentielles},
Inventiones Mathematicae \textbf{148} (2002), 219--284.
Theorems~0.6 and 5.19; the additive logarithm and its
Frobenius scaling are in \S2.4, pp.~23--24, and \S3.1,
p.~30, of the author preprint.
\url{https://doi.org/10.1007/s002220100202}.
\url{https://arxiv.org/pdf/math/0102179v3}.

\bibitem{yuan-zhang-adelic}
X. Yuan and S.-W. Zhang,
\emph{Adelic Line Bundles on Quasi-projective Varieties},
arXiv:2105.13587v9 (2025).
Theorems~A.4.2 and A.5.2; relative semipositivity and the
height normalization are in Appendix~A.
\url{https://arxiv.org/abs/2105.13587v9}.

\bibitem{genestier-ngo-shimura}
A. Genestier and B.-C. Ng\^o, \emph{Lectures on Shimura varieties},
July 2006 lecture notes, Theorem 2.2.3, p.\ 13.
\url{https://www.math.uchicago.edu/~ngo/Shimura.pdf}.

\bibitem{gaudron-remond-minima}
\'E. Gaudron and G. R\'emond, Minima, pentes et alg\`ebre tensorielle,
\emph{Israel Journal of Mathematics} \textbf{195} (2013), 565--591,
Sections 2.1--2.3, Theorem 3.1 and Lemma 3.2.
\url{https://doi.org/10.1007/s11856-012-0109-x};
\url{https://lmbp.uca.fr/~gaudron/art10.pdf}.
\bibitem{pazuki-theta-faltings}
F.~Pazuki,
\emph{Theta height and Faltings height},
Bulletin de la Soci\'et\'e Math\'ematique de France
\textbf{140} (2012), 19--49.
\url{https://www.numdam.org/item/BSMF_2012__140_1_19_0.pdf}.

\bibitem{conrad-semistable}
B.~Conrad,
\emph{Semistable reduction for abelian varieties},
lecture notes, 2011, Theorem~4.4 and Corollary~4.5, pp.~10--11.
\url{https://math.stanford.edu/~conrad/DarmonCM/2011Notes/SemistableReduction.pdf}.

\bibitem{autissier-matrix}
P.~Autissier,
\emph{Un lemme matriciel effectif},
Mathematische Zeitschrift \textbf{273} (2013), 355--361,
Corollary~1.4.
\url{https://arxiv.org/pdf/1104.5644v1}.
Author's revision:
\url{https://www.math.u-bordeaux.fr/~pautissi/Matrix2.pdf}.

\bibitem{mazur-messing-universal}
B.~Mazur and W.~Messing,
\emph{Universal Extensions and One Dimensional Crystalline Cohomology},
Lecture Notes in Mathematics 370, Springer, 1974.
Rigidified extensions: I, \S2.1, p.~13, and Proposition 2.6.7,
pp.~24--25; crystals and comparison: II, \S9, pp.~93--94,
\S13, pp.~119--122, and \S15, pp.~126--132.
\url{https://www.sas.rochester.edu/mth/sites/doug-ravenel/otherpapers/mazur-messing.pdf}.

\bibitem{conrad-units-product}
B.~Conrad, \emph{Units on product varieties},
undated notes, 4~pp., Theorem~1.1, p.~1, with proof in
\S\S2--3, pp.~1--3.
\url{https://math.stanford.edu/~conrad/papers/unitthm.pdf}.

\bibitem{madapusi-pera-realizations}
K.~Madapusi Pera,
\emph{Toroidal compactifications of integral models of Shimura
varieties of Hodge type},
Annales scientifiques de l'\'Ecole normale sup\'erieure
(4) 52 (2019), no.~2, 393--514.
For the integral one-motive comparison and the Hodge kernel of
Frobenius, see \S1.1.4, pp.~404--405, especially (1.1.4.1).
\url{https://www.numdam.org/item/10.24033/asens.2391.pdf}.

\bibitem{stacks-finite-place}
The Stacks Project Authors,
\emph{The Stacks Project},
Tags 0200 (Chow's lemma), 0AYQ (degrees on curves),
0B40 (sections of line bundles on proper curves),
and 0BW0 (the different and trace duality).
\url{https://stacks.math.columbia.edu/tag/0200};
\url{https://stacks.math.columbia.edu/tag/0AYQ};
\url{https://stacks.math.columbia.edu/tag/0B40};
\url{https://stacks.math.columbia.edu/tag/0BW0}.

\bibitem{cm-milne-av}
J.~S.~Milne, \emph{Abelian Varieties}, version 2.0, 16 March 2008.
Theorems 2.8--2.9, Remark 8.12, Proposition 13.6 and Corollary 16.4,
pp.~14--15, 39, 58--59 and 68.
\url{https://jmilne.org/math/CourseNotes/AV.pdf}.

\bibitem{cm-milne-mf}
J.~S.~Milne, \emph{Modular Functions and Modular Forms},
version 1.31, 22 March 2017.
Theorems 4.21, 12.14 and 12.18 and Corollary 12.9,
pp.~57, 124--126.
\url{https://www.jmilne.org/math/CourseNotes/MF.pdf}.

\bibitem{cm-zagier}
D.~Zagier, \emph{A Kronecker limit formula for real quadratic fields},
Mathematische Annalen \textbf{213} (1975), 153--184.
Section 2, pp.~156--159, especially (2.3)--(2.5).
\url{https://people.mpim-bonn.mpg.de/zagier/files/doi/10.1007/BF01343950/KroneckerLimit.pdf}.

\bibitem{cm-conrad-minimal}
B.~Conrad, \emph{Minimal models for elliptic curves},
21 November 2015, Definition 5.1 and Example 5.3, p.~13.
\url{https://math.stanford.edu/~conrad/papers/minimalmodel.pdf}.

\bibitem{cm-serre-tate}
J.-P.~Serre and J.~Tate, \emph{Good reduction of abelian varieties},
Annals of Mathematics (2) \textbf{88} (1968), 492--517.
Theorem 1 and Corollary 2, p.~493.
\url{https://doi.org/10.2307/1970722}.
\url{https://wstein.org/papers/bib/Serre-Tate-Good_Reduction_of_Abelian_Varieties.pdf}.

\bibitem{SerreGAGA}
J.-P.~Serre, \emph{Géométrie algébrique et géométrie analytique},
Annales de l'Institut Fourier \textbf{6} (1956), 1--42.
\S12, Theorems 1--2, pp.~19--20.
\url{https://numdam.org/articles/10.5802/aif.59/}.

\bibitem{WlodarczykResolution}
J.~W{\l}odarczyk, \emph{Simple Hironaka resolution in characteristic zero},
Journal of the American Mathematical Society \textbf{18} (2005), 779--822.
Theorems 1.0.1 and 1.0.3, pp.~781--782.
\url{https://www.math.purdue.edu/~wlodarcz/singularities/Resolution.pdf}.

\end{thebibliography}
